\begin{abstract}
We treat Daniell-style integration systems, introduce the category
$\mathbf{Int}$ whose morphisms are pullbacks along maps of the base spaces,
and locate the Daniell extension as a structure of this category.

The main result is that the Daniell extension defines a functor
$\mathcal{D}:\mathbf{Int}\to\mathbf{DInt}$ and is left adjoint to the inclusion
of the full subcategory $\mathbf{DInt}$ of Daniell closed integration systems.
That is, $\mathbf{DInt}$ is a reflective subcategory of $\mathbf{Int}$ and the
Daniell extension is its reflector.

The paper contributes the following four points.

\begin{enumerate}
  \item \textbf{Decomposing Daniell closedness into two closure conditions.}
    An integration system is Daniell closed (that is, $L^1=L$) if and only if it
    satisfies the monotone convergence theorem and is closed under modification
    on null sets. Neither condition can be dropped. The two conditions on the
    right occur independently in measure theory; the second one corresponds to
    completeness of the measure.
  \item \textbf{Determining the second reflection.}
    Requiring Daniell closedness of the codomain is not a dispensable
    restriction: the Daniell extension itself is \emph{not} left adjoint to the
    inclusion of the category of integration systems satisfying the monotone
    convergence theorem. That inclusion does have a left adjoint, but its
    reflector is \emph{not} the Daniell extension. Under the Stone condition
    this second reflection is determined as the set of all elements of $L^1$
    that are measurable with respect to the $\sigma$-algebra generated by $L$.
    This characterizes, in terms of $(X,L,I)$ itself, when the two reflections
    agree. For the integration system of step functions it consists of the
    Borel measurable integrable functions. Moreover, for an arbitrary
    integration system, $L^1$ splits as the sum of its $\sigma(L)$-measurable
    part and the functions supported on null sets. Hence \textbf{the gap between
    the two reflections amounts exactly to the functions supported on null
    sets}, which in the case of step functions appears as the difference between
    Borel and Lebesgue, that is, as the completion of the measure. Neither
    separation of points by $L$ nor the Stone condition is sufficient for the
    two reflections to agree.
  \item \textbf{Assembling the whole into a category.}
    We define identities and composition, verify that the Daniell extension
    satisfies the functor laws, and make the uniqueness of the factorization
    explicit. This part is formal. Analytic content enters the uniqueness only
    in the category obtained by dividing morphisms by null functions and in the
    category of operator morphisms below.
  \item \textbf{Making explicit what the category of operator morphisms sees.}
    We detach morphisms from maps of the base spaces and redefine them as
    positive linear maps between Riesz spaces. Operator morphisms extend
    uniquely to a codomain whose quotient is a Banach space, and the Daniell
    extension again gives a reflection in this category. However, morphisms of
    this category see only the quotient $\widetilde{L}=L/\mathcal{N}_I$: the
    assignment $(X,L,I)\mapsto(\widetilde{L},\widetilde{I})$ is faithful and
    full, so neither the base space nor the realization as a function space is
    distinguished. What the Daniell extension gives in this category is the
    Banach completion. The subcategory to be taken as the codomain of the main
    theorem is the class of objects whose quotient is a Banach space, not the
    class of Daniell closed objects; the latter is not isomorphism-closed, and
    the former is its isomorphism closure. Consequently the distinction between
    an abstract completion and a concrete realization is retained as categorical
    structure only on the side where morphisms are maps of the base spaces.
\end{enumerate}

An appendix treats the general theory of adjoint functors used in the proofs,
and the idea of forgetting the base space.

The construction of the Daniell extension by series representations, the fact
that the extension is again an integration system, that it satisfies the
monotone convergence theorem, that it is idempotent, the minimality of the
extension for a fixed base space, and the notion corresponding to a morphism of
integration systems are all known and are not claimed as new here. Among the
four points above, (1) and (2) are the parts with analytic content, but both lie
in the neighbourhood of classical Daniell theory, and the claim is only that,
\textbf{to the author's knowledge, they have not been stated explicitly}.
Point (3) is formal and point (4) is an observation about the choice of category.

Taken together, (1) and (2) split the Daniell extension into a stage that takes
the $\sigma(L)$-measurable integrable functions and a stage that completes the
measure. Both are classical operations. The position of this paper is therefore
not a new completion theorem, but rather
\textbf{the Daniell counterpart of the completion of a measure space, organized
in categorical language}.
Section~\ref{chap:intro} discusses the relation to earlier work in detail.
\end{abstract}
